\documentclass[type = bachelor]{whu-thesis}
\usepackage{array}
% \documentclass[type = master,class = academic]{whu-thesis}
% \documentclass[type = doctor]{whu-thesis}
% type: 学位类型，可选项�?bachelor, master, doctor
% class: 学位类别，可选项�?academic, professional
% showframe: 显示页面布局框架

% 以下仅列举了部分可能用到的设置选项，更多用法请参考文档《whu-thesis.pdf�?

% \PassOptionsToPackage{gbnamefmt = lowercase}{biblatex} % 英文作者姓名不强制大写

\whusetup{
  info = {
    title      = {基于有限元与机器学习的船用螺旋桨模态频率预测研究}, % 标题，可使用 \\ 手动换行
    title*     = {Research on modal frequency prediction of marine propellers based on finite element analysis and machine learning},
    department  = {弘毅学堂},
    department* = {Hongyi Honor College},
    author     = {陈子义},
    author*    = {Ziyi Chen},
    student-id = {2022302131156},
    supervisor  = {冯朝},
    supervisor* = {Chao Feng},
    academic-title  = {副研究员},
    academic-title* = {Associate Research Fellow},
    % supervisor-outer = {王某某}, % 校外导师（非必填）
    % academic-title-outer = {高级工程师}, % 校外导师职称（非必填）
    subject = {}, % 学科名称（非必填）
    major   = {智能制造工程},
    major*  = {Intelligent Manufacturing Engineering},
    research-area  = {船用螺旋桨结构振动与智能预测},
    research-area* = {Structural vibration and intelligent prediction of marine propellers},
    year = 2026,
    month = 5,
    % clc = , % 分类号
    % udc = ,
    keywords  = {船用螺旋桨, 湿模态频率预测, 高精度仿真, 模态试验, 流固耦合, 数据驱动方法, 机器学习}, % 中文关键词
    keywords* = {marine propeller, wet modal frequency prediction, high-fidelity simulation, modal test, fluid-structure coupling, data-driven method, machine learning}, % 英文关键词
  },
  style = {
    % 字体相关选项
    font = termes, % 西文字体，可选项�?default, times, xits, termes
    math-font = termes, % 数学字体，可选项�?default, xits, termes
    cjk-font = windows, % 中文字体，可选项�?windows, mac, fandol(Linux/Overleaf/TexPage), sourcehan, none
    % cjk-fakefont = true, % 使用伪粗体与伪斜�?
    % 参考文献及引用相关选项
    bib-backend = bibtex, % 参考文献引擎，可选项�?bibtex, biblatex
    % bib-style = numerical, % 参考文献样式，可选项�?numerical, author-year
    % cite-style = <>, % 引用样式（自定义�?
    bib-resource = {ref/bachelor-refs.bib}, % 参考文献数据源
    % 页面相关选项
    % chapter-page-header = true, % 章节首页是否有页�?
    % bachelor-encover = true, % 本科毕业论文英文封面
    % library, % 图书馆模式（去掉论文中所有的空白页）
    % license, % 使用授权协议�?
    % fullwidth-stop = true, % 句号样式
    % footnote-style = <>, % 脚注编号样式
    % abstract-keywords-type  = blankline, % 摘要与关键词之间样式，可选项�?blankline, newline, vfill
    % abstract-keywords-type* = blankline, % 摘要与关键词之间样式，可选项�?blankline, newline, vfill
  }
}
\whumodule{algorithm2e}
\ExplSyntaxOn
\cs_set_protected:Npn \__whu_bachelor_copyright: {}
\ExplSyntaxOff
\begin{document}
% \raggedbottom % 使得空白都置于每一页底部，可参�?https://github.com/whutug/whu-thesis/issues/276

\tableofcontents % 目录
% \listoffigures % 图目�?
% \listoftables % 表目�?

% 符号�?
% \begin{notation}
%   $\omega_n$ & $n$-维欧氏空间中单位球的表面�?\\
%   $\alpha_n$ & $n$-维欧氏空间中单位球的体积 \\
% \end{notation}

\mainmatter

\chapter{绪论}

\section{研究背景与意义}

本节内容待补充。

\section{螺旋桨模态分析研究现状}

本节内容待补充。

\section{本文研究内容与技术路线}

本节内容待补充。

\chapter{螺旋桨模态高精度仿真计算方法和计算规程}

\section{引言}

本节内容待补充。

\section{模态分析理论基础}

本节内容待补充。

\section{模态分析有限元求解理论}

本节内容待补充。

\section{螺旋桨模态分析计算规程}

本节内容待补充。

\section{本章小结}

本节内容待补充。

\chapter{螺旋桨模态高可靠测试方法和结果分析}

\section{引言}

本节内容待补充。

\section{模态测试理论基础}

本节内容待补充。

\section{螺旋桨高阶湿模态测试难点分析}

本节内容待补充。

\section{螺旋桨湿模态高可靠测试方法和测试系统}

本节内容待补充。

\section{某螺旋桨湿模态结果分析}

本节内容待补充。

\section{本章小结}

本节内容待补充。

\chapter{螺旋桨湿模态预报高性能机器学习方法}

\section{引言}

本节内容待补充。

\section{机器学习基本理论}

本节内容待补充。

\section{螺旋桨湿模态机器学习预报难点分析}

本节内容待补充。

\section{螺旋桨模态预报高性能预报流程}

本节内容待补充。

\section{螺旋桨湿模态机器学习预报结果分析}

本节内容待补充。

\section{本章小结}

本节内容待补充。

\chapter{总结与展望}

\section{全文总结}

本文面向船用螺旋桨模态频率快速预测问题，围绕“有限元高精度仿真--模态测试验证--机器学习快速预测”这一技术路线开展研究。通过建立螺旋桨干/湿模态有限元仿真规程、开展模态测试验证，并构建干模态频率数据库，训练机器学习预测模型，最终形成了基于干模态预测结果和附连水修正系数的湿模态频率快速预报方法。

在螺旋桨模态仿真方面，本文在有限元模态分析理论的基础上，建立了螺旋桨干模态与湿模态仿真计算规程。针对螺旋桨几何形状复杂、叶片厚度变化明显、局部模态对网格质量敏感等特点，从网格类型、网格尺寸、材料参数设置和求解流程等方面进行了规范化分析，确定了适用于螺旋桨干模态计算的有限元建模方法。在湿模态仿真方面，建立了结构域与流体域耦合模型，通过声学流体单元和结构--流体界面约束实现湿模态频率求解。以 DTMB4382 型螺旋桨为对象，应用该规程开展干/湿模态计算，获得了前 20 阶干模态频率和湿模态频率，并计算得到前 20 阶附连水频率修正系数，为后续湿模态快速预测提供了基础。

在数据库构建方面，本文基于参数化建模和自动化仿真技术，实现了螺旋桨干模态频率数据库的自动化构建。通过 Python 与 CadQuery 生成螺旋桨三维实体模型，利用 HyperMesh 批处理脚本完成模型导入与网格自动划分，并借助 Abaqus Python 脚本实现模态分析任务自动提交、结果读取和数据汇总。该流程将几何建模、网格划分、有限元求解和频率提取等步骤串联起来，形成了较完整的自动化仿真流程，构建了干模态频率数据库，为后续机器学习模型训练提供了数据基础。

在螺旋桨模态测试方面，本文在讲述模态测试基本原理的基础上，开展了螺旋桨干模态与湿模态测试研究。针对高阶湿模态测试中存在的振型复杂、模态成簇、频响峰值衰减和信噪比较低等问题，建立了基于移动力锤法、加速度响应频响函数测试和模态参数辨识的测试流程。分别以某 10 叶螺旋桨和某 7 叶螺旋桨为对象，开展空气中和水中条件下的模态测试，获得了不同螺旋桨在干/湿模态下的频率与振型特征。其中，对 7 叶螺旋桨进一步开展干/湿模态有限元计算，并将仿真结果与测试结果进行对比。结果表明，计算频率与测试频率吻合较好，主要振型对应关系明确，说明本文建立的有限元仿真规程能够较准确地反映螺旋桨干/湿模态动力学特性。

在机器学习预测方面，本文基于干模态频率数据库开展模型训练。首先采用 Pearson 相关系数和 Spearman 相关系数分析原始几何参数与模态频率之间的关系，筛选出相关系数均大于 0.1 的 24 个几何参数作为模型输入，实现输入变量降维。随后，分别训练随机森林回归、梯度提升回归、高斯过程回归和多层感知机四种模型，并采用随机搜索算法进行超参数优化。利用 MAE、RMSE、$R^2$ 和 MAPE 四种指标对模型性能进行评价后发现，多层感知机模型表现最优，其测试集整体 MAE 为 5.133937，RMSE 为 7.445019，$R^2$ 为 0.996914，MAPE 为 1.889931\%。该结果说明 MLP 能够较好学习螺旋桨关键几何参数与前 20 阶干模态频率之间的非线性映射关系，预测结果与有限元计算结果吻合程度较高。最后，本文将 MLP 预测得到的干模态频率与附连水频率修正系数结合，提出了“几何参数输入--干模态频率预测--附连水系数修正--湿模态频率输出”的螺旋桨湿模态频率快速预测流程。

\section{未来展望}

本文围绕船用螺旋桨模态频率快速预测问题，建立了基于有限元仿真与机器学习代理模型的预测方法，并实现了由几何参数到干、湿模态频率的快速估计。但受研究周期、样本规模和建模方法的限制，仍存在进一步完善的空间：

可以进一步开展高阶模态预测研究。本文主要针对螺旋桨前 20 阶模态频率开展预测研究，预测对象以低阶和中阶模态为主，而对 20 阶以上高阶局部模态的研究仍不充分。实际工作中，高阶局部模态可能与叶片局部振动、噪声辐射以及复杂流场激励有关，后续可进一步构建包含更高阶模态频率的数据库，提升模型对高阶局部振动特性的预测能力。

可以进一步提高湿模态频率预测精度。本文湿模态频率预测采用“干模态频率预测结果结合附连水系数修正”的方法。该方法计算效率较高，但附连水系数本身受到螺旋桨几何形状、模态阶次、流体域建模范围和流固耦合计算精度等因素影响，因此其修正精度仍存在一定局限。后续可进一步扩大湿模态仿真样本数量，直接建立“螺旋桨几何参数--湿模态频率”数据库，从而减少附连水系数近似带来的误差。

可以进一步扩展螺旋桨几何变量范围。本文样本集中螺旋桨翼型相对单一，主要关注直径、桨叶数、弦长和厚度等几何参数对模态频率的影响，尚未将翼型变化作为重要设计变量纳入模型输入。后续研究可将不同翼型参数引入样本设计，进一步扩充数据库规模，使机器学习模型能够学习翼型差异对螺旋桨刚度分布、质量分布及模态频率的影响规律，从而提高预测模型的适用范围和工程推广价值。



% 参考文�?
% \nocite{*}
\nocite{intro:1,intro:2,intro:3,intro:4,intro:5,intro:6,intro:7,intro:8,intro:9,intro:10,intro:11,intro:12,intro:13,intro:14,intro:15,intro:16,intro:17,intro:18,intro:19,sim:20,sim:21}
\printbibliography

% 发表的与学位论文相关的科研成果目�?
% \include{pages/achievement.tex}

% 致谢
\begin{acknowledgements}
  致谢是以简短的文字对课题研究与论文撰写过程中直接给予帮助的人员(例如指导教师、答疑教师及其他人员)表示谢意。致谢是对他人劳动的尊重，也是学术规范。内容限一页。
\end{acknowledgements}

% 附录
\end{document}
